% Options for packages loaded elsewhere
\PassOptionsToPackage{unicode}{hyperref}
\PassOptionsToPackage{hyphens}{url}
\PassOptionsToPackage{dvipsnames,svgnames,x11names}{xcolor}
%
\documentclass[
  letterpaper,
  DIV=11,
  numbers=noendperiod]{scrartcl}

\usepackage{amsmath,amssymb}
\usepackage{iftex}
\ifPDFTeX
  \usepackage[T1]{fontenc}
  \usepackage[utf8]{inputenc}
  \usepackage{textcomp} % provide euro and other symbols
\else % if luatex or xetex
  \usepackage{unicode-math}
  \defaultfontfeatures{Scale=MatchLowercase}
  \defaultfontfeatures[\rmfamily]{Ligatures=TeX,Scale=1}
\fi
\usepackage{lmodern}
\ifPDFTeX\else  
    % xetex/luatex font selection
\fi
% Use upquote if available, for straight quotes in verbatim environments
\IfFileExists{upquote.sty}{\usepackage{upquote}}{}
\IfFileExists{microtype.sty}{% use microtype if available
  \usepackage[]{microtype}
  \UseMicrotypeSet[protrusion]{basicmath} % disable protrusion for tt fonts
}{}
\makeatletter
\@ifundefined{KOMAClassName}{% if non-KOMA class
  \IfFileExists{parskip.sty}{%
    \usepackage{parskip}
  }{% else
    \setlength{\parindent}{0pt}
    \setlength{\parskip}{6pt plus 2pt minus 1pt}}
}{% if KOMA class
  \KOMAoptions{parskip=half}}
\makeatother
\usepackage{xcolor}
\setlength{\emergencystretch}{3em} % prevent overfull lines
\setcounter{secnumdepth}{-\maxdimen} % remove section numbering
% Make \paragraph and \subparagraph free-standing
\makeatletter
\ifx\paragraph\undefined\else
  \let\oldparagraph\paragraph
  \renewcommand{\paragraph}{
    \@ifstar
      \xxxParagraphStar
      \xxxParagraphNoStar
  }
  \newcommand{\xxxParagraphStar}[1]{\oldparagraph*{#1}\mbox{}}
  \newcommand{\xxxParagraphNoStar}[1]{\oldparagraph{#1}\mbox{}}
\fi
\ifx\subparagraph\undefined\else
  \let\oldsubparagraph\subparagraph
  \renewcommand{\subparagraph}{
    \@ifstar
      \xxxSubParagraphStar
      \xxxSubParagraphNoStar
  }
  \newcommand{\xxxSubParagraphStar}[1]{\oldsubparagraph*{#1}\mbox{}}
  \newcommand{\xxxSubParagraphNoStar}[1]{\oldsubparagraph{#1}\mbox{}}
\fi
\makeatother

\usepackage{color}
\usepackage{fancyvrb}
\newcommand{\VerbBar}{|}
\newcommand{\VERB}{\Verb[commandchars=\\\{\}]}
\DefineVerbatimEnvironment{Highlighting}{Verbatim}{commandchars=\\\{\}}
% Add ',fontsize=\small' for more characters per line
\usepackage{framed}
\definecolor{shadecolor}{RGB}{241,243,245}
\newenvironment{Shaded}{\begin{snugshade}}{\end{snugshade}}
\newcommand{\AlertTok}[1]{\textcolor[rgb]{0.68,0.00,0.00}{#1}}
\newcommand{\AnnotationTok}[1]{\textcolor[rgb]{0.37,0.37,0.37}{#1}}
\newcommand{\AttributeTok}[1]{\textcolor[rgb]{0.40,0.45,0.13}{#1}}
\newcommand{\BaseNTok}[1]{\textcolor[rgb]{0.68,0.00,0.00}{#1}}
\newcommand{\BuiltInTok}[1]{\textcolor[rgb]{0.00,0.23,0.31}{#1}}
\newcommand{\CharTok}[1]{\textcolor[rgb]{0.13,0.47,0.30}{#1}}
\newcommand{\CommentTok}[1]{\textcolor[rgb]{0.37,0.37,0.37}{#1}}
\newcommand{\CommentVarTok}[1]{\textcolor[rgb]{0.37,0.37,0.37}{\textit{#1}}}
\newcommand{\ConstantTok}[1]{\textcolor[rgb]{0.56,0.35,0.01}{#1}}
\newcommand{\ControlFlowTok}[1]{\textcolor[rgb]{0.00,0.23,0.31}{\textbf{#1}}}
\newcommand{\DataTypeTok}[1]{\textcolor[rgb]{0.68,0.00,0.00}{#1}}
\newcommand{\DecValTok}[1]{\textcolor[rgb]{0.68,0.00,0.00}{#1}}
\newcommand{\DocumentationTok}[1]{\textcolor[rgb]{0.37,0.37,0.37}{\textit{#1}}}
\newcommand{\ErrorTok}[1]{\textcolor[rgb]{0.68,0.00,0.00}{#1}}
\newcommand{\ExtensionTok}[1]{\textcolor[rgb]{0.00,0.23,0.31}{#1}}
\newcommand{\FloatTok}[1]{\textcolor[rgb]{0.68,0.00,0.00}{#1}}
\newcommand{\FunctionTok}[1]{\textcolor[rgb]{0.28,0.35,0.67}{#1}}
\newcommand{\ImportTok}[1]{\textcolor[rgb]{0.00,0.46,0.62}{#1}}
\newcommand{\InformationTok}[1]{\textcolor[rgb]{0.37,0.37,0.37}{#1}}
\newcommand{\KeywordTok}[1]{\textcolor[rgb]{0.00,0.23,0.31}{\textbf{#1}}}
\newcommand{\NormalTok}[1]{\textcolor[rgb]{0.00,0.23,0.31}{#1}}
\newcommand{\OperatorTok}[1]{\textcolor[rgb]{0.37,0.37,0.37}{#1}}
\newcommand{\OtherTok}[1]{\textcolor[rgb]{0.00,0.23,0.31}{#1}}
\newcommand{\PreprocessorTok}[1]{\textcolor[rgb]{0.68,0.00,0.00}{#1}}
\newcommand{\RegionMarkerTok}[1]{\textcolor[rgb]{0.00,0.23,0.31}{#1}}
\newcommand{\SpecialCharTok}[1]{\textcolor[rgb]{0.37,0.37,0.37}{#1}}
\newcommand{\SpecialStringTok}[1]{\textcolor[rgb]{0.13,0.47,0.30}{#1}}
\newcommand{\StringTok}[1]{\textcolor[rgb]{0.13,0.47,0.30}{#1}}
\newcommand{\VariableTok}[1]{\textcolor[rgb]{0.07,0.07,0.07}{#1}}
\newcommand{\VerbatimStringTok}[1]{\textcolor[rgb]{0.13,0.47,0.30}{#1}}
\newcommand{\WarningTok}[1]{\textcolor[rgb]{0.37,0.37,0.37}{\textit{#1}}}

\providecommand{\tightlist}{%
  \setlength{\itemsep}{0pt}\setlength{\parskip}{0pt}}\usepackage{longtable,booktabs,array}
\usepackage{calc} % for calculating minipage widths
% Correct order of tables after \paragraph or \subparagraph
\usepackage{etoolbox}
\makeatletter
\patchcmd\longtable{\par}{\if@noskipsec\mbox{}\fi\par}{}{}
\makeatother
% Allow footnotes in longtable head/foot
\IfFileExists{footnotehyper.sty}{\usepackage{footnotehyper}}{\usepackage{footnote}}
\makesavenoteenv{longtable}
\usepackage{graphicx}
\makeatletter
\def\maxwidth{\ifdim\Gin@nat@width>\linewidth\linewidth\else\Gin@nat@width\fi}
\def\maxheight{\ifdim\Gin@nat@height>\textheight\textheight\else\Gin@nat@height\fi}
\makeatother
% Scale images if necessary, so that they will not overflow the page
% margins by default, and it is still possible to overwrite the defaults
% using explicit options in \includegraphics[width, height, ...]{}
\setkeys{Gin}{width=\maxwidth,height=\maxheight,keepaspectratio}
% Set default figure placement to htbp
\makeatletter
\def\fps@figure{htbp}
\makeatother

\KOMAoption{captions}{tableheading}
\makeatletter
\@ifpackageloaded{caption}{}{\usepackage{caption}}
\AtBeginDocument{%
\ifdefined\contentsname
  \renewcommand*\contentsname{Table of contents}
\else
  \newcommand\contentsname{Table of contents}
\fi
\ifdefined\listfigurename
  \renewcommand*\listfigurename{List of Figures}
\else
  \newcommand\listfigurename{List of Figures}
\fi
\ifdefined\listtablename
  \renewcommand*\listtablename{List of Tables}
\else
  \newcommand\listtablename{List of Tables}
\fi
\ifdefined\figurename
  \renewcommand*\figurename{Figure}
\else
  \newcommand\figurename{Figure}
\fi
\ifdefined\tablename
  \renewcommand*\tablename{Table}
\else
  \newcommand\tablename{Table}
\fi
}
\@ifpackageloaded{float}{}{\usepackage{float}}
\floatstyle{ruled}
\@ifundefined{c@chapter}{\newfloat{codelisting}{h}{lop}}{\newfloat{codelisting}{h}{lop}[chapter]}
\floatname{codelisting}{Listing}
\newcommand*\listoflistings{\listof{codelisting}{List of Listings}}
\makeatother
\makeatletter
\makeatother
\makeatletter
\@ifpackageloaded{caption}{}{\usepackage{caption}}
\@ifpackageloaded{subcaption}{}{\usepackage{subcaption}}
\makeatother

\ifLuaTeX
  \usepackage{selnolig}  % disable illegal ligatures
\fi
\usepackage{bookmark}

\IfFileExists{xurl.sty}{\usepackage{xurl}}{} % add URL line breaks if available
\urlstyle{same} % disable monospaced font for URLs
\hypersetup{
  pdftitle={Evaluating Predictive Model Accuracy in Foreclosure Cases},
  pdfauthor={Piru Puiggari},
  colorlinks=true,
  linkcolor={blue},
  filecolor={Maroon},
  citecolor={Blue},
  urlcolor={Blue},
  pdfcreator={LaTeX via pandoc}}


\title{Evaluating Predictive Model Accuracy in Foreclosure Cases}
\author{Piru Puiggari}
\date{2024-08-23}

\begin{document}
\maketitle


\subsection{Explanation}

\subsubsection{In this tab:}\label{in-this-tab}

\begin{itemize}
\tightlist
\item
  \hyperref[sec-introduction]{Introduction}
\item
  \hyperref[sec-dataset]{Dataset Description}
\item
  \hyperref[sec-methodology]{Methodology}
\item
  \hyperref[sec-current-model]{Current Model Approach}
\item
  \hyperref[sec-random-forest]{Random Forest Model}
\item
  \hyperref[sec-recommendations]{Recommendations}
\end{itemize}

\subsubsection{Introduction}\label{sec-introduction}

Accurate prediction of foreclosure timelines is critical for optimizing
the valuation of distressed mortgage assets. Foreclosure dates play a
crucial role in estimating asset recovery values within a mortgage
portfolio, serving as a baseline for portfolio valuation and the
structuring of settlements or other termination strategies. The ability
to forecast these timelines with precision allows for better risk
management, enhanced asset recovery, and more informed decisions
regarding portfolio performance.

This report examines the current predictive model for foreclosure times
and introduces an enhanced approach aimed at improving the accuracy and
reliability of these predictions. The new model focuses on predicting
the duration required to resolve each legal proceeding within
foreclosure cases, offering significant improvements in accuracy. These
advancements are essential for supporting strategic financial decisions
and effectively managing Non-Performing Loans (NPLs), which often
involve high Loan-to-Value (LTV) ratios and extended delinquency periods

\subsubsection{Dataset Description}\label{sec-dataset}

The dataset used in this project contains detailed information on over
1,200 mortgage foreclosure cases, each of which has been fully resolved.
For each case, the dataset includes all the legal proceedings that
occurred throughout its duration. This comprehensive data is crucial for
understanding the factors that influence the time it takes to resolve
these cases, providing a solid foundation for the predictive models
developed in this project. Key columns in the dataset include:

\begin{itemize}
\item
  \textbf{Portfolio Name}: Identifies the portfolio to which the legal
  case belongs.
\item
  \textbf{Legal File Name}: A unique identifier for each legal case.
\item
  \textbf{Date}: The date of each legal proceeding.
\item
  \textbf{Legal Stage}: The stage of the legal process
\item
  \textbf{Legal Proceeding}: The specific type of legal action taken
  within the case.
\item
  \textbf{Foreclosure Initial Date}: The date when the foreclosure
  process began.
\item
  \textbf{Resolution Date}: The date when the case was resolved.
\item
  \textbf{Time to Solve}: The duration between the date of the
  proceeding and the resolution of the case.
\item
  \textbf{Geographical Information}: Including \textbf{Zone},
  \textbf{Region}, and \textbf{Zone Group}, which provide context on
  where the case is being handled.
\item
  \textbf{Case Characteristics}: Such as whether an appeal (cassation)
  was involved (\textbf{Tiene Casacion}) and the central stage of the
  legal process (\textbf{Etapa Central}).
\end{itemize}

\subsubsection{Methodology for Evaluating Regression
Models}\label{sec-methodology}

To ensure our predictive model's accuracy and reliability, we use three
key evaluation metrics:

\begin{enumerate}
\def\labelenumi{\arabic{enumi}.}
\tightlist
\item
  \textbf{Mean Absolute Error (MAE)}

  \begin{itemize}
  \tightlist
  \item
    Measures the average difference between predicted and actual
    resolution times
  \item
    Interpreted as: On average, how many months off are the model's
    predictions?
  \item
    Lower MAE indicates higher accuracy
  \end{itemize}
\item
  \textbf{Root Mean Squared Error (RMSE)}

  \begin{itemize}
  \tightlist
  \item
    Similar to MAE, but gives more weight to larger errors
  \item
    Provides insight into the magnitude of the model's prediction
    mistakes
  \item
    Lower RMSE suggests better model performance
  \end{itemize}
\item
  \textbf{R-squared (R²)}

  \begin{itemize}
  \tightlist
  \item
    Explains how well the model's predictions match actual outcomes
  \item
    Expressed as a percentage (0-100\%)
  \item
    Higher R² indicates the model captures more factors influencing
    resolution times
  \item
    Example: R² of 0.94 means the model explains 94\% of variations in
    resolution times
  \end{itemize}
\end{enumerate}

\subsubsection{Approach for Current Model}\label{sec-current-model}

The current model predicts the resolution time for foreclosure cases by
using historical averages. Specifically, it looks at past cases and
calculates the average time it took to resolve cases based on the type
of legal proceeding involved. For each new case, the model predicts the
resolution time by matching it to similar cases from the past and
applying the average time for those cases. The process involves:

\begin{enumerate}
\def\labelenumi{\arabic{enumi}.}
\item
  \textbf{Data Collection \& Cleaning:} We gathered and cleaned past
  foreclosure data to ensure reliability.
\item
  \textbf{Grouping by Legal Proceeding:} Cases were categorized by legal
  proceeding type, and average resolution times were calculated for each
  group.
\item
  \textbf{Prediction:} The model predicts resolution time based on the
  type of proceeding, using historical averages. If no match is found,
  an overall average is used. This straightforward approach assumes that
  past performance predicts future outcomes.
\end{enumerate}

\paragraph{Evaluation of the Current Model's
Performance}\label{evaluation-of-the-current-models-performance}

To evaluate the effectiveness of our current model, we tested it against
foreclosure cases that have already been resolved. This allows us to
compare the model's predictions with the actual resolution times and
measure its accuracy. The key metrics we used are:

\begin{itemize}
\item
  \textbf{Mean Absolute Error (MAE):} The model's average prediction was
  off by 16.68 months. This means that, on average, the model's
  predictions are almost 17 months away from the actual resolution time.
\item
  \textbf{Root Mean Squared Error (RMSE):} This metric, which emphasizes
  larger errors, shows a typical prediction error of 22.06 months. The
  higher error indicates that the model struggles with cases that
  deviate significantly from the average.
\item
  \textbf{R-squared (R²):} The model explains only 39\% of the variance
  in resolution times. This low percentage suggests that the model does
  not capture many of the factors influencing how long a case will take.
\end{itemize}

While the model provides some predictive capability, its accuracy and
reliability are inadequate. The significant errors and low variance
explanation suggest the need for a more sophisticated approach.

\paragraph{Analysis of the Actual vs.~Predicted Graph for the Current
Model}\label{analysis-of-the-actual-vs.-predicted-graph-for-the-current-model}

The Actual vs.~Predicted graph compares the model's predictions with the
actual resolution times for foreclosure cases. In an ideal scenario, all
data points would lie along the diagonal line (red dashed line),
indicating that the predicted values perfectly match the actual values.

\begin{figure}[H]

{\centering \includegraphics{actual vs.png}

}

\caption{Actual vs.~Predicted Resolution Times for Current Model}

\end{figure}%

However, in the graph for the current model:

\begin{itemize}
\item
  \textbf{Wide Dispersion:} Many points are scattered far from the
  diagonal line, indicating large prediction errors. This spread
  reflects the model's inaccuracies, particularly in cases where the
  actual resolution time deviates from the average.
\item
  \textbf{Horizontal and Vertical Lines:} These occur because the model
  uses average times for specific legal proceedings, leading to
  identical predictions for different cases. This approach doesn't
  account for the variability within those cases, resulting in clusters
  of predictions that don't match the actual outcomes.
\end{itemize}

For example, in a specific case where the actual resolution time was 30
months, the model might have predicted only 15 months based on the
average for that type of proceeding. This discrepancy highlights the
model's tendency to underestimate or overestimate the resolution time,
depending on how closely a case aligns with the historical average.

\paragraph{Limitations of Using Averages for Predicting Time to
Solve}\label{limitations-of-using-averages-for-predicting-time-to-solve}

The graphs in this section displays the distribution of resolution times
for the 20 most common legal proceedings. Each bar or line in the graph
represents how resolution times are spread out for these specific types
of proceedings.

\begin{figure}[H]

{\centering \includegraphics{distribuciones.png}

}

\caption{Distribution of Resolution Times for Top 20 Legal Proceedings}

\end{figure}%

\begin{itemize}
\item
  \textbf{Wide Variance:} For many legal proceedings, there is a broad
  range of resolution times. When the distribution of these times is far
  from the mean (average), using that average as a prediction leads to
  significant errors. This is because the average fails to capture the
  true spread of the data, especially in cases with outliers or skewed
  distributions.
\item
  \textbf{Skewed Distributions:} In some cases, the distribution is
  heavily skewed, meaning that a few cases take much longer or shorter
  than most others. The average gets pulled in one direction, making it
  a poor predictor for the majority of cases.
\end{itemize}

\begin{figure}[H]

{\centering \includegraphics{output.png}

}

\caption{Feature Importance for Random Forest Model}

\end{figure}%

\begin{itemize}
\tightlist
\item
  \textbf{Non-Normal Distributions:} Many proceedings have skewed or
  uneven distributions, meaning that most cases are not near the average
  time. Using a simple average as a predictor fails to account for these
  nuances, leading to significant errors.
\end{itemize}

In essence, relying on averages oversimplifies the prediction process
and fails to account for the variability in actual case outcomes. This
can lead to substantial inaccuracies, especially when the distribution
is not centered around the mean

\subsubsection{Random Forest Predictive Model}\label{sec-random-forest}

To enhance the accuracy of our foreclosure time predictions, we
developed a Random Forest model. This advanced machine learning
technique constructs multiple decision trees, each one analyzing
different aspects of the data. By averaging the predictions from these
trees, the model effectively manages the complexities of the dataset and
minimizes the risk of overfitting.

Our model leverages a wide range of features, such as the sequence of
legal proceedings, time intervals between actions, and specific
attributes related to the courts handling the cases. By considering
these diverse factors, the Random Forest model delivers predictions that
are significantly more accurate and reliable

\paragraph{Model Development}\label{model-development}

1.~~~~ \textbf{Data Preprocessing:} Cleaned dataset, removed irrelevant
columns, handled missing values, and filtered for foreclosure cases.
This ensures we're working with accurate, relevant data.

2.~~~~ \textbf{Feature Engineering:} We created features such as
Months\_Since\_Start, Proceeding\_Count,
Months\_Since\_Last\_Proceeding, Average\_Months\_Between\_Proceedings,
Distinct\_Proceeding\_Types, and Repeated\_Proceeding\_Count. These
features capture key temporal, procedural, and case-specific aspects,
improving prediction accuracy.

3.~~~~ \textbf{Encoding and Scaling}: Applied Target Encoding for
high-cardinality variables, Label Encoding for other categorical
features, and StandardScaler for numerical features. This prepares the
data for optimal model performance.

4.~~~~ \textbf{Model Training:} Split the dataset, then trained a Random
Forest with 300 trees. This approach leverages multiple decision trees
to make robust predictions about foreclosure resolution times.

\paragraph{Results}\label{results}

The Random Forest model's performance in predicting foreclosure
resolution times is summarized by the following key metrics:

\begin{itemize}
\item
  \textbf{Mean Absolute Error (MAE):} 4.89 months
\item
  \textbf{Root Mean Squared Error (RMSE):} 6.99 months
\item
  \textbf{R-squared (R²):} 0.94
\end{itemize}

These metrics indicate a high level of accuracy, with the model
explaining 94\% of the variance in the resolution times, and relatively
low prediction errors.

\begin{figure}[H]

{\centering \includegraphics{rf pred.png}

}

\caption{Random Forest Model Predictions vs Actual Resolution Times}

\end{figure}%

\paragraph{Feature Importance}\label{feature-importance}

The bar chart ranks the importance of various features used by the model
to make predictions. Added Features like Months\_Since\_Start, and
Proceeding\_Count emerge as the most significant factors influencing the
resolution time.

Understanding which features most strongly affect the model's
predictions can provide valuable insights. For instance, the high
importance of Months\_Since\_Start and Proceeding\_Count highlights the
relevance of the duration and frequency of legal actions in determining
how long a foreclosure case will take to resolve.

\begin{figure}[H]

{\centering \includegraphics{rf importance.png}

}

\caption{Feature Importance for Random Forest Model}

\end{figure}%

This visualization provides a clear ranking of the features that have
the most significant impact on the model's predictions. It offers
valuable insights into which factors are most influential in determining
foreclosure resolution times.

\subsubsection{Recommendations}\label{sec-recommendations}

\begin{enumerate}
\def\labelenumi{\arabic{enumi}.}
\item
  \textbf{Expand Features:} \textbf{Expand Features:} Enhance the
  model's accuracy by incorporating additional features, such as
  Loan-to-Value (LTV) ratios to help predict whether a foreclosure will
  be resolved by a third party or the creditor, along with economic
  indicators or more detailed geographic data.
\item
  \textbf{Explore Advanced Techniques:} Consider testing more advanced
  machine learning techniques in the future for further improvements
\item
  \textbf{Integrate with Financial Tools:} Integrate the model's
  predictions into existing financial systems to support seamless
  decision-making and help structure DPO (Discounted Payoff) agreements
  more effectively.
\end{enumerate}

\subsection{Code \& Analysis}

\subsubsection{In this section:}\label{in-this-section}

\begin{itemize}
\tightlist
\item
  \hyperref[sec-data-loading]{Data Loading}
\item
  \hyperref[sec-feature-engineering]{Feature Engineering}
\item
  \hyperref[sec-encoding]{Encoding and Preparation}
\item
  \hyperref[sec-model-training]{Model Training}
\item
  \hyperref[sec-visualization]{Visualization}
\item
  \hyperref[sec-conclusion]{Conclusion}
\end{itemize}

\subsubsection{Data Loading and Preprocessing}\label{sec-data-loading}

First, we'll import the necessary libraries and load our dataset. We're
working with an Excel file containing information about mortgage
foreclosure cases.

\begin{Shaded}
\begin{Highlighting}[]
\ImportTok{import}\NormalTok{ pandas }\ImportTok{as}\NormalTok{ pd}
\ImportTok{import}\NormalTok{ numpy }\ImportTok{as}\NormalTok{ np}
\ImportTok{import}\NormalTok{ matplotlib.pyplot }\ImportTok{as}\NormalTok{ plt}
\ImportTok{import}\NormalTok{ seaborn }\ImportTok{as}\NormalTok{ sns}
\ImportTok{from}\NormalTok{ sklearn.model\_selection }\ImportTok{import}\NormalTok{ train\_test\_split}
\ImportTok{from}\NormalTok{ sklearn.preprocessing }\ImportTok{import}\NormalTok{ StandardScaler, LabelEncoder}
\ImportTok{from}\NormalTok{ sklearn.ensemble }\ImportTok{import}\NormalTok{ RandomForestRegressor}
\ImportTok{from}\NormalTok{ sklearn.metrics }\ImportTok{import}\NormalTok{ mean\_absolute\_error, mean\_squared\_error, r2\_score}
\ImportTok{from}\NormalTok{ category\_encoders }\ImportTok{import}\NormalTok{ TargetEncoder}

\CommentTok{\# Load the data}
\NormalTok{file\_path }\OperatorTok{=} \StringTok{"/Users/piru/Desktop/tiempos actuaciones v2 por juzgados v3.xlsx"}
\NormalTok{sheet\_name }\OperatorTok{=} \StringTok{"Base3"}
\NormalTok{df }\OperatorTok{=}\NormalTok{ pd.read\_excel(file\_path, sheet\_name}\OperatorTok{=}\NormalTok{sheet\_name)}
\end{Highlighting}
\end{Shaded}

Now that we have our data loaded, we need to clean it. This involves
handling missing values, converting date fields, and ensuring all
categorical variables are in the correct format.

\begin{Shaded}
\begin{Highlighting}[]
\CommentTok{\# Clean and preprocess data}
\NormalTok{df }\OperatorTok{=}\NormalTok{ df[}\OperatorTok{\textasciitilde{}}\NormalTok{df[}\StringTok{\textquotesingle{}foreclosure initial date\textquotesingle{}}\NormalTok{].}\BuiltInTok{apply}\NormalTok{(}\KeywordTok{lambda}\NormalTok{ x: }\BuiltInTok{isinstance}\NormalTok{(x, time))]}
\NormalTok{df.dropna(axis}\OperatorTok{=}\DecValTok{1}\NormalTok{, how}\OperatorTok{=}\StringTok{\textquotesingle{}all\textquotesingle{}}\NormalTok{, inplace}\OperatorTok{=}\VariableTok{True}\NormalTok{)}
\NormalTok{df.dropna(axis}\OperatorTok{=}\DecValTok{0}\NormalTok{, how}\OperatorTok{=}\StringTok{\textquotesingle{}all\textquotesingle{}}\NormalTok{, inplace}\OperatorTok{=}\VariableTok{True}\NormalTok{)}
\NormalTok{df }\OperatorTok{=}\NormalTok{ df.dropna(subset}\OperatorTok{=}\NormalTok{[}\StringTok{\textquotesingle{}Legal Stage\textquotesingle{}}\NormalTok{])}

\CommentTok{\# Convert dates and clean text fields}
\NormalTok{df[}\StringTok{\textquotesingle{}foreclosure initial date\textquotesingle{}}\NormalTok{] }\OperatorTok{=}\NormalTok{ pd.to\_datetime(df[}\StringTok{\textquotesingle{}foreclosure initial date\textquotesingle{}}\NormalTok{])}
\NormalTok{df[}\StringTok{\textquotesingle{}resolution date\textquotesingle{}}\NormalTok{] }\OperatorTok{=}\NormalTok{ pd.to\_datetime(df[}\StringTok{\textquotesingle{}resolution date\textquotesingle{}}\NormalTok{], errors}\OperatorTok{=}\StringTok{\textquotesingle{}coerce\textquotesingle{}}\NormalTok{)}
\NormalTok{df }\OperatorTok{=}\NormalTok{ df.}\BuiltInTok{apply}\NormalTok{(}\KeywordTok{lambda}\NormalTok{ x: x.}\BuiltInTok{str}\NormalTok{.upper() }\ControlFlowTok{if}\NormalTok{ x.dtype }\OperatorTok{==} \StringTok{\textquotesingle{}object\textquotesingle{}} \ControlFlowTok{else}\NormalTok{ x)}

\CommentTok{\# Convert boolean columns}
\ControlFlowTok{if} \StringTok{\textquotesingle{}etapa central\textquotesingle{}} \KeywordTok{in}\NormalTok{ df.columns:}
\NormalTok{    df[}\StringTok{\textquotesingle{}etapa central\textquotesingle{}}\NormalTok{] }\OperatorTok{=}\NormalTok{ df[}\StringTok{\textquotesingle{}etapa central\textquotesingle{}}\NormalTok{].astype(}\StringTok{\textquotesingle{}bool\textquotesingle{}}\NormalTok{)}
\ControlFlowTok{if} \StringTok{\textquotesingle{}tiene casacion\textquotesingle{}} \KeywordTok{in}\NormalTok{ df.columns:}
\NormalTok{    df[}\StringTok{\textquotesingle{}tiene casacion\textquotesingle{}}\NormalTok{] }\OperatorTok{=}\NormalTok{ df[}\StringTok{\textquotesingle{}tiene casacion\textquotesingle{}}\NormalTok{].astype(}\StringTok{\textquotesingle{}bool\textquotesingle{}}\NormalTok{)}
\end{Highlighting}
\end{Shaded}

We also need to ensure that our categorical variables are properly
typed. This is important for later encoding steps and for the Random
Forest model to handle these variables correctly.

\begin{Shaded}
\begin{Highlighting}[]
\CommentTok{\# Define and apply data types}
\NormalTok{data\_types }\OperatorTok{=}\NormalTok{ \{}
    \StringTok{\textquotesingle{}Legal Stage\textquotesingle{}}\NormalTok{: }\StringTok{\textquotesingle{}category\textquotesingle{}}\NormalTok{,}
    \StringTok{\textquotesingle{}Portfolio Name\textquotesingle{}}\NormalTok{: }\StringTok{\textquotesingle{}category\textquotesingle{}}\NormalTok{,}
    \StringTok{\textquotesingle{}Legal Proceeding\textquotesingle{}}\NormalTok{: }\StringTok{\textquotesingle{}category\textquotesingle{}}\NormalTok{,}
    \StringTok{\textquotesingle{}clave\textquotesingle{}}\NormalTok{: }\StringTok{\textquotesingle{}category\textquotesingle{}}\NormalTok{,}
    \StringTok{\textquotesingle{}resolution type\textquotesingle{}}\NormalTok{: }\StringTok{\textquotesingle{}category\textquotesingle{}}\NormalTok{,}
    \StringTok{\textquotesingle{}general type\textquotesingle{}}\NormalTok{: }\StringTok{\textquotesingle{}category\textquotesingle{}}\NormalTok{,}
    \StringTok{\textquotesingle{}zone\textquotesingle{}}\NormalTok{: }\StringTok{\textquotesingle{}category\textquotesingle{}}\NormalTok{,}
    \StringTok{\textquotesingle{}zone group\textquotesingle{}}\NormalTok{: }\StringTok{\textquotesingle{}category\textquotesingle{}}\NormalTok{,}
    \StringTok{\textquotesingle{}region\textquotesingle{}}\NormalTok{: }\StringTok{\textquotesingle{}category\textquotesingle{}}\NormalTok{,}
    \StringTok{\textquotesingle{}Law Firm\textquotesingle{}}\NormalTok{: }\StringTok{\textquotesingle{}category\textquotesingle{}}\NormalTok{,}
    \StringTok{\textquotesingle{}Judicial Court\textquotesingle{}}\NormalTok{: }\StringTok{\textquotesingle{}category\textquotesingle{}}\NormalTok{,}
\NormalTok{\}}

\ControlFlowTok{for}\NormalTok{ col, dtype }\KeywordTok{in}\NormalTok{ data\_types.items():}
    \ControlFlowTok{if}\NormalTok{ col }\KeywordTok{in}\NormalTok{ df.columns:}
\NormalTok{        df[col] }\OperatorTok{=}\NormalTok{ df[col].astype(dtype)}
\end{Highlighting}
\end{Shaded}

\subsubsection{Feature Engineering}\label{sec-feature-engineering}

Now that our data is clean, we can start creating new features that
might help predict the resolution time. We'll focus on temporal aspects
of the legal proceedings.

First, let's filter our data to include only foreclosure cases with a
positive resolution time:

\begin{Shaded}
\begin{Highlighting}[]
\CommentTok{\# Filter dataframe for FORECLOSURE cases with time to solve \textgreater{} 0.01}
\NormalTok{df }\OperatorTok{=}\NormalTok{ df[(df[}\StringTok{\textquotesingle{}general type\textquotesingle{}}\NormalTok{] }\OperatorTok{==} \StringTok{\textquotesingle{}FORECLOSURE\textquotesingle{}}\NormalTok{) }\OperatorTok{\&}\NormalTok{ (df[}\StringTok{\textquotesingle{}time to solve\textquotesingle{}}\NormalTok{] }\OperatorTok{\textgreater{}} \FloatTok{0.01}\NormalTok{)]}

\CommentTok{\# Remove legal proceedings with less than 5 occurrences}
\NormalTok{legal\_proceeding\_counts }\OperatorTok{=}\NormalTok{ df[}\StringTok{\textquotesingle{}Legal Proceeding\textquotesingle{}}\NormalTok{].value\_counts()}
\NormalTok{valid\_proceedings }\OperatorTok{=}\NormalTok{ legal\_proceeding\_counts[legal\_proceeding\_counts }\OperatorTok{\textgreater{}=} \DecValTok{5}\NormalTok{].index}
\NormalTok{df }\OperatorTok{=}\NormalTok{ df[df[}\StringTok{\textquotesingle{}Legal Proceeding\textquotesingle{}}\NormalTok{].isin(valid\_proceedings)]}

\CommentTok{\# Ensure the dataframe is sorted by Legal File Name and Date}
\NormalTok{df }\OperatorTok{=}\NormalTok{ df.sort\_values([}\StringTok{\textquotesingle{}Legal File Name\textquotesingle{}}\NormalTok{, }\StringTok{\textquotesingle{}Date\textquotesingle{}}\NormalTok{])}
\end{Highlighting}
\end{Shaded}

Now, let's create our new features:

\begin{Shaded}
\begin{Highlighting}[]
\CommentTok{\# Calculate new features}
\NormalTok{df[}\StringTok{\textquotesingle{}Case\_Start\textquotesingle{}}\NormalTok{] }\OperatorTok{=}\NormalTok{ df.groupby(}\StringTok{\textquotesingle{}Legal File Name\textquotesingle{}}\NormalTok{)[}\StringTok{\textquotesingle{}Date\textquotesingle{}}\NormalTok{].transform(}\StringTok{\textquotesingle{}min\textquotesingle{}}\NormalTok{)}
\NormalTok{df[}\StringTok{\textquotesingle{}Months\_Since\_Start\textquotesingle{}}\NormalTok{] }\OperatorTok{=}\NormalTok{ (df[}\StringTok{\textquotesingle{}Date\textquotesingle{}}\NormalTok{] }\OperatorTok{{-}}\NormalTok{ df[}\StringTok{\textquotesingle{}Case\_Start\textquotesingle{}}\NormalTok{]).dt.days }\OperatorTok{/} \FloatTok{30.44}
\NormalTok{df[}\StringTok{\textquotesingle{}Proceeding\_Count\textquotesingle{}}\NormalTok{] }\OperatorTok{=}\NormalTok{ df.groupby(}\StringTok{\textquotesingle{}Legal File Name\textquotesingle{}}\NormalTok{).cumcount() }\OperatorTok{+} \DecValTok{1}
\NormalTok{df[}\StringTok{\textquotesingle{}Months\_Since\_Last\_Proceeding\textquotesingle{}}\NormalTok{] }\OperatorTok{=}\NormalTok{ df.groupby(}\StringTok{\textquotesingle{}Legal File Name\textquotesingle{}}\NormalTok{)[}\StringTok{\textquotesingle{}Date\textquotesingle{}}\NormalTok{].diff().dt.days }\OperatorTok{/} \FloatTok{30.44}
\NormalTok{df[}\StringTok{\textquotesingle{}Months\_Since\_Last\_Proceeding\textquotesingle{}}\NormalTok{] }\OperatorTok{=}\NormalTok{ df[}\StringTok{\textquotesingle{}Months\_Since\_Last\_Proceeding\textquotesingle{}}\NormalTok{].fillna(}\DecValTok{1}\NormalTok{)}
\NormalTok{df[}\StringTok{\textquotesingle{}Average\_Months\_Between\_Proceedings\textquotesingle{}}\NormalTok{] }\OperatorTok{=}\NormalTok{ df.groupby(}\StringTok{\textquotesingle{}Legal File Name\textquotesingle{}}\NormalTok{)[}\StringTok{\textquotesingle{}Months\_Since\_Last\_Proceeding\textquotesingle{}}\NormalTok{].transform(}\StringTok{\textquotesingle{}mean\textquotesingle{}}\NormalTok{)}
\NormalTok{df[}\StringTok{\textquotesingle{}Distinct\_Proceeding\_Types\textquotesingle{}}\NormalTok{] }\OperatorTok{=}\NormalTok{ df.groupby(}\StringTok{\textquotesingle{}Legal File Name\textquotesingle{}}\NormalTok{)[}\StringTok{\textquotesingle{}Legal Proceeding\textquotesingle{}}\NormalTok{].transform(}\StringTok{\textquotesingle{}nunique\textquotesingle{}}\NormalTok{)}
\NormalTok{df[}\StringTok{\textquotesingle{}Repeated\_Proceeding\_Count\textquotesingle{}}\NormalTok{] }\OperatorTok{=}\NormalTok{ df.groupby([}\StringTok{\textquotesingle{}Legal File Name\textquotesingle{}}\NormalTok{, }\StringTok{\textquotesingle{}Legal Proceeding\textquotesingle{}}\NormalTok{]).cumcount() }\OperatorTok{+} \DecValTok{1}
\NormalTok{df[}\StringTok{\textquotesingle{}Previous\_Proceeding\textquotesingle{}}\NormalTok{] }\OperatorTok{=}\NormalTok{ df.groupby(}\StringTok{\textquotesingle{}Legal File Name\textquotesingle{}}\NormalTok{)[}\StringTok{\textquotesingle{}Legal Proceeding\textquotesingle{}}\NormalTok{].shift(}\DecValTok{1}\NormalTok{)}
\NormalTok{df[}\StringTok{\textquotesingle{}Previous\_Proceeding\textquotesingle{}}\NormalTok{] }\OperatorTok{=}\NormalTok{ df[}\StringTok{\textquotesingle{}Previous\_Proceeding\textquotesingle{}}\NormalTok{].fillna(df[}\StringTok{\textquotesingle{}Legal Proceeding\textquotesingle{}}\NormalTok{])}

\CommentTok{\# Handle missing zone group}
\NormalTok{df[}\StringTok{\textquotesingle{}zone group\textquotesingle{}}\NormalTok{] }\OperatorTok{=}\NormalTok{ df[}\StringTok{\textquotesingle{}zone group\textquotesingle{}}\NormalTok{].fillna(df[}\StringTok{\textquotesingle{}region\textquotesingle{}}\NormalTok{])}
\NormalTok{df[}\StringTok{\textquotesingle{}zone group\textquotesingle{}}\NormalTok{] }\OperatorTok{=}\NormalTok{ pd.Categorical(df[}\StringTok{\textquotesingle{}zone group\textquotesingle{}}\NormalTok{])}
\end{Highlighting}
\end{Shaded}

These new features capture various temporal aspects of each case, such
as how long the case has been ongoing, how many proceedings have
occurred, and the frequency of proceedings.

\subsubsection{Encoding and Model Preparation}\label{sec-encoding}

Before we can train our model, we need to encode our categorical
variables. We'll use Target Encoding for high-cardinality variables and
Label Encoding for others.

\begin{Shaded}
\begin{Highlighting}[]
\CommentTok{\# Target Encoding for high{-}cardinality categorical variables}
\NormalTok{high\_cardinality\_cols }\OperatorTok{=}\NormalTok{ [}\StringTok{\textquotesingle{}Legal Proceeding\textquotesingle{}}\NormalTok{, }\StringTok{\textquotesingle{}Judicial Court\textquotesingle{}}\NormalTok{, }\StringTok{\textquotesingle{}Law Firm\textquotesingle{}}\NormalTok{, }\StringTok{\textquotesingle{}clave\textquotesingle{}}\NormalTok{, }\StringTok{\textquotesingle{}Previous\_Proceeding\textquotesingle{}}\NormalTok{]}
\NormalTok{te }\OperatorTok{=}\NormalTok{ TargetEncoder()}
\NormalTok{df\_encoded }\OperatorTok{=}\NormalTok{ te.fit\_transform(df[high\_cardinality\_cols], df[}\StringTok{\textquotesingle{}time to solve\textquotesingle{}}\NormalTok{])}

\CommentTok{\# Replace original columns with target encoded versions}
\ControlFlowTok{for}\NormalTok{ col }\KeywordTok{in}\NormalTok{ high\_cardinality\_cols:}
\NormalTok{    df[col] }\OperatorTok{=}\NormalTok{ df\_encoded[col]}

\CommentTok{\# Label Encoding for other categorical variables}
\NormalTok{le }\OperatorTok{=}\NormalTok{ LabelEncoder()}
\NormalTok{categorical\_columns\_to\_encode }\OperatorTok{=}\NormalTok{ [}\StringTok{\textquotesingle{}Legal Stage\textquotesingle{}}\NormalTok{, }\StringTok{\textquotesingle{}Portfolio Name\textquotesingle{}}\NormalTok{, }\StringTok{\textquotesingle{}region\textquotesingle{}}\NormalTok{, }\StringTok{\textquotesingle{}zone group\textquotesingle{}}\NormalTok{]}
\ControlFlowTok{for}\NormalTok{ col }\KeywordTok{in}\NormalTok{ categorical\_columns\_to\_encode:}
\NormalTok{    df[col] }\OperatorTok{=}\NormalTok{ le.fit\_transform(df[col].astype(}\BuiltInTok{str}\NormalTok{))}
\NormalTok{    df[col] }\OperatorTok{=}\NormalTok{ pd.Categorical(df[col])}
\end{Highlighting}
\end{Shaded}

Now we can define our features and prepare our data for modeling:

\begin{Shaded}
\begin{Highlighting}[]
\CommentTok{\# Define features for modeling}
\NormalTok{features }\OperatorTok{=}\NormalTok{ [}
    \StringTok{\textquotesingle{}Legal Stage\textquotesingle{}}\NormalTok{, }\StringTok{\textquotesingle{}Legal Proceeding\textquotesingle{}}\NormalTok{, }\StringTok{\textquotesingle{}Portfolio Name\textquotesingle{}}\NormalTok{, }\StringTok{\textquotesingle{}zone group\textquotesingle{}}\NormalTok{, }\StringTok{\textquotesingle{}region\textquotesingle{}}\NormalTok{,}
    \StringTok{\textquotesingle{}Judicial Court\textquotesingle{}}\NormalTok{, }\StringTok{\textquotesingle{}Law Firm\textquotesingle{}}\NormalTok{, }\StringTok{\textquotesingle{}clave\textquotesingle{}}\NormalTok{, }\StringTok{\textquotesingle{}etapa central\textquotesingle{}}\NormalTok{, }\StringTok{\textquotesingle{}Months\_Since\_Start\textquotesingle{}}\NormalTok{,}
    \StringTok{\textquotesingle{}Proceeding\_Count\textquotesingle{}}\NormalTok{, }\StringTok{\textquotesingle{}Months\_Since\_Last\_Proceeding\textquotesingle{}}\NormalTok{, }\StringTok{\textquotesingle{}Average\_Months\_Between\_Proceedings\textquotesingle{}}\NormalTok{,}
    \StringTok{\textquotesingle{}Distinct\_Proceeding\_Types\textquotesingle{}}\NormalTok{, }\StringTok{\textquotesingle{}Repeated\_Proceeding\_Count\textquotesingle{}}\NormalTok{, }\StringTok{\textquotesingle{}Previous\_Proceeding\textquotesingle{}}\NormalTok{,}
    \StringTok{\textquotesingle{}time to solve\textquotesingle{}}
\NormalTok{]}

\CommentTok{\# Prepare X and y for modeling}
\NormalTok{target\_col }\OperatorTok{=} \StringTok{\textquotesingle{}time to solve\textquotesingle{}}
\NormalTok{X }\OperatorTok{=}\NormalTok{ df[features].drop(columns}\OperatorTok{=}\NormalTok{[target\_col])}
\NormalTok{y }\OperatorTok{=}\NormalTok{ df[target\_col]}
\end{Highlighting}
\end{Shaded}

\subsubsection{Model Training and Evaluation}\label{sec-model-training}

We'll use a Random Forest Regressor for our prediction task. First,
let's split our data into training and testing sets:

\begin{Shaded}
\begin{Highlighting}[]
\CommentTok{\# Split the data}
\NormalTok{X\_train, X\_test, y\_train, y\_test }\OperatorTok{=}\NormalTok{ train\_test\_split(X, y, test\_size}\OperatorTok{=}\FloatTok{0.2}\NormalTok{, random\_state}\OperatorTok{=}\DecValTok{42}\NormalTok{)}
\end{Highlighting}
\end{Shaded}

Next, we need to scale our numerical features:

\begin{Shaded}
\begin{Highlighting}[]
\CommentTok{\# Scale numerical features}
\NormalTok{numerical\_columns }\OperatorTok{=}\NormalTok{ X.select\_dtypes(include}\OperatorTok{=}\NormalTok{[}\StringTok{\textquotesingle{}int64\textquotesingle{}}\NormalTok{, }\StringTok{\textquotesingle{}float64\textquotesingle{}}\NormalTok{]).columns}
\NormalTok{scaler }\OperatorTok{=}\NormalTok{ StandardScaler()}
\NormalTok{X\_train\_scaled }\OperatorTok{=}\NormalTok{ X\_train.copy()}
\NormalTok{X\_test\_scaled }\OperatorTok{=}\NormalTok{ X\_test.copy()}
\NormalTok{X\_train\_scaled[numerical\_columns] }\OperatorTok{=}\NormalTok{ scaler.fit\_transform(X\_train[numerical\_columns])}
\NormalTok{X\_test\_scaled[numerical\_columns] }\OperatorTok{=}\NormalTok{ scaler.transform(X\_test[numerical\_columns])}
\end{Highlighting}
\end{Shaded}

Now we can train our Random Forest model:

\begin{Shaded}
\begin{Highlighting}[]
\CommentTok{\# Train Random Forest model}
\NormalTok{rf\_model }\OperatorTok{=}\NormalTok{ RandomForestRegressor(n\_estimators}\OperatorTok{=}\DecValTok{300}\NormalTok{, min\_samples\_split}\OperatorTok{=}\DecValTok{2}\NormalTok{, }
\NormalTok{                                 min\_samples\_leaf}\OperatorTok{=}\DecValTok{1}\NormalTok{, max\_depth}\OperatorTok{=}\VariableTok{None}\NormalTok{, }
\NormalTok{                                 bootstrap}\OperatorTok{=}\VariableTok{True}\NormalTok{, random\_state}\OperatorTok{=}\DecValTok{42}\NormalTok{)}
\NormalTok{rf\_model.fit(X\_train\_scaled, y\_train)}
\end{Highlighting}
\end{Shaded}

Finally, let's make predictions and evaluate our model:

\begin{Shaded}
\begin{Highlighting}[]
\CommentTok{\# Make predictions and evaluate model}
\NormalTok{y\_pred }\OperatorTok{=}\NormalTok{ rf\_model.predict(X\_test\_scaled)}
\NormalTok{mae }\OperatorTok{=}\NormalTok{ mean\_absolute\_error(y\_test, y\_pred)}
\NormalTok{mse }\OperatorTok{=}\NormalTok{ mean\_squared\_error(y\_test, y\_pred)}
\NormalTok{rmse }\OperatorTok{=}\NormalTok{ np.sqrt(mse)}
\NormalTok{r2 }\OperatorTok{=}\NormalTok{ r2\_score(y\_test, y\_pred)}

\BuiltInTok{print}\NormalTok{(}\SpecialStringTok{f"Mean Absolute Error (MAE): }\SpecialCharTok{\{}\NormalTok{mae}\SpecialCharTok{\}}\SpecialStringTok{"}\NormalTok{)}
\BuiltInTok{print}\NormalTok{(}\SpecialStringTok{f"Root Mean Squared Error (RMSE): }\SpecialCharTok{\{}\NormalTok{rmse}\SpecialCharTok{\}}\SpecialStringTok{"}\NormalTok{)}
\BuiltInTok{print}\NormalTok{(}\SpecialStringTok{f"R{-}squared (R2): }\SpecialCharTok{\{}\NormalTok{r2}\SpecialCharTok{\}}\SpecialStringTok{"}\NormalTok{)}
\end{Highlighting}
\end{Shaded}

Mean Absolute Error (MAE): 4.887682724118453

Root Mean Squared Error (RMSE): 6.985917898599015

R-squared (R2): 0.9382800622825924

\subsubsection{Visualization}\label{sec-visualization}

To better understand our model's performance and the importance of our
features, we'll create three visualizations.

First, let's look at the feature importance:

\begin{Shaded}
\begin{Highlighting}[]
\CommentTok{\# Feature Importance plot}
\NormalTok{feature\_importance }\OperatorTok{=}\NormalTok{ pd.DataFrame(\{}
    \StringTok{\textquotesingle{}feature\textquotesingle{}}\NormalTok{: X.columns,}
    \StringTok{\textquotesingle{}importance\textquotesingle{}}\NormalTok{: rf\_model.feature\_importances\_}
\NormalTok{\}).sort\_values(}\StringTok{\textquotesingle{}importance\textquotesingle{}}\NormalTok{, ascending}\OperatorTok{=}\VariableTok{False}\NormalTok{)}

\NormalTok{plt.figure(figsize}\OperatorTok{=}\NormalTok{(}\DecValTok{12}\NormalTok{, }\DecValTok{8}\NormalTok{))}
\NormalTok{sns.barplot(x}\OperatorTok{=}\StringTok{\textquotesingle{}importance\textquotesingle{}}\NormalTok{, y}\OperatorTok{=}\StringTok{\textquotesingle{}feature\textquotesingle{}}\NormalTok{, data}\OperatorTok{=}\NormalTok{feature\_importance.head(}\DecValTok{20}\NormalTok{))}
\NormalTok{plt.title(}\StringTok{\textquotesingle{}Top 20 Feature Importance from Random Forest\textquotesingle{}}\NormalTok{)}
\NormalTok{plt.tight\_layout()}
\NormalTok{plt.show()}
\end{Highlighting}
\end{Shaded}

\begin{figure}[H]

{\centering \includegraphics{rf importance.png}

}

\caption{Feature Importance for Random Forest Model}

\end{figure}%

This plot shows us which features are most influential in our model's
predictions.

Next, let's compare our predictions to the actual values:

\begin{Shaded}
\begin{Highlighting}[]
\CommentTok{\# Actual vs Predicted plot}
\NormalTok{plt.figure(figsize}\OperatorTok{=}\NormalTok{(}\DecValTok{10}\NormalTok{, }\DecValTok{6}\NormalTok{))}
\NormalTok{plt.scatter(y\_test, y\_pred, alpha}\OperatorTok{=}\FloatTok{0.5}\NormalTok{)}
\NormalTok{plt.plot([y\_test.}\BuiltInTok{min}\NormalTok{(), y\_test.}\BuiltInTok{max}\NormalTok{()], [y\_test.}\BuiltInTok{min}\NormalTok{(), y\_test.}\BuiltInTok{max}\NormalTok{()], }\StringTok{\textquotesingle{}r{-}{-}\textquotesingle{}}\NormalTok{, lw}\OperatorTok{=}\DecValTok{2}\NormalTok{)}
\NormalTok{plt.xlabel(}\StringTok{\textquotesingle{}Actual\textquotesingle{}}\NormalTok{)}
\NormalTok{plt.ylabel(}\StringTok{\textquotesingle{}Predicted\textquotesingle{}}\NormalTok{)}
\NormalTok{plt.title(}\StringTok{\textquotesingle{}Actual vs Predicted\textquotesingle{}}\NormalTok{)}
\NormalTok{plt.tight\_layout()}
\NormalTok{plt.show()}
\end{Highlighting}
\end{Shaded}

\includegraphics{rf pred.png} This plot helps us visualize how well our
predictions match the actual values. Points closer to the red line
indicate better predictions.

Finally, let's look at the residuals of our model:

\begin{Shaded}
\begin{Highlighting}[]
\CommentTok{\# Residual plot}
\NormalTok{residuals }\OperatorTok{=}\NormalTok{ y\_test }\OperatorTok{{-}}\NormalTok{ y\_pred}
\NormalTok{plt.figure(figsize}\OperatorTok{=}\NormalTok{(}\DecValTok{10}\NormalTok{, }\DecValTok{6}\NormalTok{))}
\NormalTok{plt.scatter(y\_pred, residuals, alpha}\OperatorTok{=}\FloatTok{0.5}\NormalTok{)}
\NormalTok{plt.axhline(y}\OperatorTok{=}\DecValTok{0}\NormalTok{, color}\OperatorTok{=}\StringTok{\textquotesingle{}r\textquotesingle{}}\NormalTok{, linestyle}\OperatorTok{=}\StringTok{\textquotesingle{}{-}{-}\textquotesingle{}}\NormalTok{)}
\NormalTok{plt.xlabel(}\StringTok{\textquotesingle{}Predicted\textquotesingle{}}\NormalTok{)}
\NormalTok{plt.ylabel(}\StringTok{\textquotesingle{}Residuals\textquotesingle{}}\NormalTok{)}
\NormalTok{plt.title(}\StringTok{\textquotesingle{}Residual Plot\textquotesingle{}}\NormalTok{)}
\NormalTok{plt.tight\_layout()}
\NormalTok{plt.show()}
\end{Highlighting}
\end{Shaded}

\begin{figure}[H]

{\centering \includegraphics{residuals.png}

}

\caption{Residual Plot for Random Forest Model}

\end{figure}%

This residual plot visualizes our model's prediction errors, offering
insights into its performance. It helps identify potential biases,
assess accuracy across different prediction ranges, and detect error
patterns. The scatter of points around the zero line (red dashed) shows
how errors vary with predicted values. A random, evenly distributed
scatter suggests a well-performing model, while clear patterns may
indicate areas where the model struggles. This visualization is key for
evaluating the model's reliability and pinpointing areas for improvement
in our forecasting approach.

\subsubsection{Conclusion}\label{conclusion}

This analysis demonstrates the use of a Random Forest model to predict
resolution times for mortgage foreclosure cases. The model shows good
performance with an R-squared value of about 0.94, indicating that it
explains a large portion of the variance in resolution times.

Key findings include:

\begin{enumerate}
\def\labelenumi{\arabic{enumi}.}
\tightlist
\item
  The importance of temporal features like `Months\_Since\_Start' and
  `Proceeding\_Count' in predicting resolution times.
\item
  The model's ability to capture both linear and non-linear
  relationships in the data.
\item
  Areas for potential improvement, as seen in the residual plot.
\end{enumerate}

Future work could involve feature selection, hyperparameter tuning, or
exploring other machine learning algorithms to potentially improve the
model's performance.




\end{document}
